% ==================================================================
% BAGIAN 1: ABSTRAK INGGRIS (HALAMAN PERTAMA)
% ==================================================================
\clearpage
\phantomsection % Penting agar hyperlink PDF akurat
\addcontentsline{toc}{chapter}{ABSTRAK} % Memasukkan 'ABSTRAK' ke Daftar Isi
\thispagestyle{plain}
\vspace*{-1cm}

\begin{center}

{\fontfamily{ptm}\selectfont\bfseries
Universitas Bina Nusantara\\[0.1cm]
\rule{\textwidth}{0.8pt}\\[0.2cm]

Computer Science Program\\
Computer Science Study Program\\
School of Computer Science\\
Skripsi Sarjana Teknik Informatika\\[0.5cm]

% --- TITLE ---
DEVELOPMENT OF AN ENVIRONMENTAL PROBLEM REPORTING SYSTEM\\
INTEGRATING IMAGE AND TEXT DATA TO SUPPORT SMART CITY\\
INITIATIVES\\[0.8cm]
}

% --- AUTHORS ---
{\fontfamily{ptm}\selectfont\bfseries
\begin{tabular}{p{7cm} l}
Brian Theodore & 2602080030 \\[0.1cm]
Moh. Khoirul Umam Al Amin & 2602198472 \\[0.1cm]
Andrea Octaviani & 2602075895
\end{tabular}
}
\end{center}

\vspace{0.3cm}

% --- HEADING ABSTRACT ---
\begin{center}
{\fontfamily{ptm}\selectfont\bfseries\itshape ABSTRACT}
\end{center}

% --- ISI ABSTRACT ---
{\fontfamily{ptm}\selectfont\itshape
\noindent The smart city concept requires the utilization of information technology to manage urban environmental issues in an effective and efficient manner. One of the challenges faced is the lack of an optimal environmental issue reporting system that can provide accurate and well-structured information. This research aims to develop an environmental issue reporting system by integrating image and text data as the main sources of report information. The developed system allows users to submit reports in the form of digital images and textual descriptions, which are then processed and stored in a centralized database. Data processing is conducted to support the classification and validation of environmental issue reports. The system is developed using the System Development Life Cycle (SDLC) method, which includes the stages of analysis, design, implementation, and testing. The testing results indicate that the system improves the quality of reporting information, accelerates data management processes, and assists relevant stakeholders in monitoring and handling environmental issues. Therefore, this system can support the implementation of a sustainable, technology-driven smart city.

\noindent\textbf{Keywords:} smart city, reporting system, image processing, text processing, informatics engineering.
}

% ==================================================================
% BAGIAN 2: ABSTRAK INDONESIA (HALAMAN KEDUA)
% ==================================================================
\newpage
\thispagestyle{plain}
\vspace*{-1cm}


\begin{center}
{\fontfamily{ptm}\selectfont\bfseries
Universitas Bina Nusantara\\[0.1cm]
\rule{\textwidth}{0.8pt}\\[0.2cm]

Computer Science Program\\
Computer Science Study Program\\
School of Computer Science\\
Skripsi Sarjana Teknik Informatika\\[0.5cm]

% --- TITLE---
PENGEMBANGAN SISTEM PELAPORAN MASALAH LINGKUNGAN DENGAN\\
INTEGRASI DATA GAMBAR DAN TEKS UNTUK MENDUKUNG \textit{SMART CITY}\\[0.8cm]
}

% --- AUTHORS ---
{\fontfamily{ptm}\selectfont\bfseries
\begin{tabular}{p{7cm} l}
Brian Theodore & 2602080030 \\[0.1cm]
Moh. Khoirul Umam Al Amin & 2602198472 \\[0.1cm]
Andrea Octaviani & 2602075895
\end{tabular}
}
\end{center}

\vspace{0.3cm}

% --- HEADING ABSTRAK ---
\begin{center}
{\fontfamily{ptm}\selectfont\bfseries ABSTRAK}
\end{center}

% --- ISI ABSTRAK ---
{\fontfamily{ptm}\selectfont
\noindent Konsep \textit{smart city} membutuhkan pemanfaatan teknologi informasi untuk mengelola masalah lingkungan perkotaan secara efektif dan efisien. Salah satu tantangan yang dihadapi adalah kurangnya sistem pelaporan masalah lingkungan yang optimal yang dapat memberikan informasi yang akurat dan terstruktur dengan baik. Penelitian ini bertujuan untuk mengembangkan sistem pelaporan masalah lingkungan dengan mengintegrasikan data gambar dan teks sebagai sumber utama informasi laporan. Sistem yang dikembangkan memungkinkan pengguna untuk mengirimkan laporan dalam bentuk gambar digital dan deskripsi teks, yang kemudian diproses dan disimpan dalam basis data terpusat. Pemrosesan data dilakukan untuk mendukung klasifikasi dan validasi laporan masalah lingkungan. Sistem dikembangkan menggunakan metode \textit{System Development Life Cycle} (SDLC), yang mencakup tahapan analisis, perancangan, implementasi, dan pengujian. Hasil pengujian menunjukkan bahwa sistem meningkatkan kualitas informasi pelaporan, mempercepat proses pengelolaan data, dan membantu pemangku kepentingan terkait dalam memantau dan menangani masalah lingkungan. Oleh karena itu, sistem ini dapat mendukung implementasi \textit{smart city} yang berkelanjutan dan berbasis teknologi.

\noindent\textbf{Kata Kunci:} \textit{smart city}, sistem pelaporan, pengolahan citra, pemrosesan teks, teknik informatika.
}

\newpage